\chapter{Weitere betrachtete Verfahren}\label{chap:weitere_verfahren}

Neben den im Hauptvergleich untersuchten Verfahren wurde im Rahmen dieser Arbeit
auch der Einsatz probabilistischer Ähnlichkeitsverfahren evaluiert. Das folgende
Kapitel beschreibt den Ansatz auf Basis von MinHash und LSH (Locality Sensitive
Hashing) sowie die Begründung, weshalb er für das kontinuierliche Log-Clustering
nicht geeignet ist.

\section{MinHash und LSH: Probabilistische Ähnlichkeitssuche}\label{sec:minhash_lsh}

MinHash und LSH bilden eine Kombination probabilistischer Algorithmen, die darauf
ausgelegt sind, ähnliche Datenpunkte in großen Mengen effizient aufzufinden. Das
Verfahren opfert dabei ein geringes Maß an Präzision zugunsten einer erheblichen
Reduktion des Rechenaufwands. Klassische Einsatzgebiete sind die
Duplikaterkennung in Dokumentensammlungen sowie kollaborative Filtersysteme im
großflächigen E-Commerce\footnote{\cite{Leskovec2020Mining} Kapitel 3.1}.

\paragraph{Funktionsprinzip.}
Das Verfahren gliedert sich in drei aufeinanderfolgende Schritte. Im ersten
Schritt wird jede Lognachricht durch \textit{Shingling} in eine mathematische
Menge überführt: Ein $k$-Shingle ist ein beliebiger Teilstring der Länge $k$,
der in der Nachricht vorkommt. Die Wahl von $k$ beeinflusst dabei maßgeblich
die Trennschärfe – zu kleine Werte führen zu fälschlich hohen
Ähnlichkeitswerten\footnote{\cite{Leskovec2020Mining} Kapitel 3.3.2}.

Als Ähnlichkeitsmetrik dient die \textit{Jaccard-Ähnlichkeit}, die das
Verhältnis von Schnitt- zu Vereinigungsmenge zweier Shingle-Mengen $S$ und $T$
beschreibt:
\begin{equation}
    \mathrm{SIM}(S, T) = \frac{|S \cap T|}{|S \cup T|}
    \label{eq:jaccard}
\end{equation}

Da der direkte mengenbasierte Vergleich bei $n$ Lognachrichten einen Aufwand
von $\mathcal{O}(n^2)$ erzeugt, reduziert \textit{MinHashing} die Mengen auf
kompakte Signaturen fester Länge $k$. Der Algorithmus permutiert die Zeilen
der booleschen Shingle-Matrix zufällig und speichert je Spalte den Index der
ersten Zeile mit dem Wert~1. Die Wahrscheinlichkeit, dass zwei Spalten nach
einer Permutation denselben MinHash-Wert aufweisen, entspricht exakt ihrer
Jaccard-Ähnlichkeit. Durch $k$-fache Wiederholung entsteht eine Signatur der
Länge $k$\footnote{\cite{Leskovec2020Mining} Kapitel 3.3}.

\textit{Locality Sensitive Hashing (LSH)} löst das verbleibende
Skalierungsproblem: Selbst mit Signaturen sind noch $\mathcal{O}(n^2)$
paarweise Vergleiche erforderlich. Das Banding-Verfahren unterteilt die
Signaturmatrix horizontal in $b$ Bänder à $r$ Zeilen und bildet jeden
Band-Teilvektor per Hashfunktion auf einen Bucket ab. Zwei Signaturen gelten
als Kandidatenpaar, sobald sie in mindestens einem Band in denselben Bucket
fallen. Die Kandidatenwahrscheinlichkeit bei Jaccard-Ähnlichkeit $s$ beträgt:
\begin{equation}
    P(\text{Kandidat}) = 1 - (1 - s^r)^b
    \label{eq:lsh_prob}
\end{equation}
Durch Variation von $b$ und $r$ lässt sich die Schwellwertkurve so steuern, dass
der Übergang von „kein Kandidat" zu „Kandidat" einem annähernden Stufensprung
gleicht\footnote{\cite{Leskovec2020Mining} Kapitel 3.4.2}.

\paragraph{Limitierungen für kontinuierliches Log-Clustering.}
Trotz der theoretischen Attraktivität des Verfahrens ergeben sich für den
vorliegenden Anwendungsfall strukturelle Einschränkungen, die einen Einsatz
im Hauptvergleich ausschließen.

Erstens ist MinHash\,+\,LSH ein \textit{Batch-Verfahren}: Die Signaturmatrix und
der LSH-Index werden über den vollständigen Datensatz aufgebaut. Bei einem
kontinuierlichen System, das periodisch neue Logs erhält, müsste der Index bei
jedem neuen Batch vollständig neu berechnet oder durch aufwendige
Merge-Operationen erweitert werden. Ein inhärentes Gedächtnis – wie es Drain3
durch seinen persistierten Parse-Baum oder der pgvector-Ansatz durch die
gespeicherten Cluster-Embeddings bereitstellt – fehlt.

Zweitens misst die Jaccard-Ähnlichkeit auf Basis von Wort-N-Grammen die
\textit{strukturelle Zeichenüberschneidung}, nicht die semantische Verwandtschaft
von Log-Templates. Variable Bestandteile wie IP-Adressen, Zeitstempel oder
numerische IDs erzeugen bei Logs desselben Templates unterschiedliche
Shingle-Mengen und senken dadurch die gemessene Ähnlichkeit künstlich ab.

\paragraph{Empirische Einordnung.}
Im Rahmen dieser Arbeit wurde das Verfahren prototypisch mit der Python-Bibliothek
\texttt{datasketch} auf dem OpenStack-Datensatz des Loghub-Benchmarks erprobt.
Bei einer wortbasierten Trigramm-Tokenisierung (\texttt{MINHASH\_NGRAM\_SIZE = 3}),
128~Permutationen (\texttt{MINHASH\_NUM\_PERM = 128}) und einem Jaccard-Schwellenwert
von $0{,}5$ (\texttt{MINHASH\_THRESHOLD = 0.5}) entstanden 1\,809~Cluster. Die
tatsächliche Anzahl der Log-Templates im Datensatz liegt jedoch deutlich darunter.
Dieses Ergebnis bestätigt die oben beschriebene Schwäche: Der gewählte
Schwellenwert reicht nicht aus, um Logs desselben Templates zuverlässig in einem
Bucket zusammenzuführen, da ihre Trigramm-Überschneidung die Schwelle bei
variablen Token-Positionen häufig unterschreitet.

\paragraph{Fazit.}
MinHash und LSH sind leistungsfähige Werkzeuge für die skalierbare
Ähnlichkeitssuche in statischen Datenmengen, erfüllen jedoch die zentralen
Anforderungen des kontinuierlichen Log-Clusterings nicht: Es fehlt sowohl ein
persistierbarer Zustand für inkrementelle Verarbeitung als auch eine
Ähnlichkeitsmetrik, die gegenüber den variablen Bestandteilen von Lognachrichten
robust ist. Das Verfahren scheidet daher als eigenständiger Clustering-Ansatz
für den vorliegenden Anwendungsfall aus.
